\documentclass[10pt,letterpaper]{article}
\usepackage{cornell-notes}

\title{Annex D.27.03: Class Template Wbuffer Convert}
\author{}
\date{}
\setDocTitle{Annex D.27.03: Class Template Wbuffer Convert}
\setDocSubtitle{C++ 2024}
\setDocOwner{C++ 2024 Cornell Notes}
\setDocDate{\today}
\setCornellCollection{C++ 2024 Cornell Notes}
\setCornellUnitType{Annex}
\setCornellUnitNumber{D.27.03}
\setCornellUnitTitle{Class Template Wbuffer Convert}
\hypersetup{pdftitle={Annex D.27.03: Class Template Wbuffer Convert},pdfsubject={Cornell notes},pdfauthor={}}

\begin{document}
\maketitle
\begin{CornellOverviewBox}{Overview}
\textbf{Classification:} Normative annex\\[2pt]
\textbf{Learning objective:} Understand the role and technical guidance associated with Class template wbuffer\_convert.
\end{CornellOverviewBox}\begin{CornellSummaryBox}{Summary}
{\sffamily\bfseries\color{Accent} Summary}\par\vspace{3pt}Section D.27.3 focuses on Class template wbuffer\_convert. Apply the specified interface together with its mandates, effects, result conditions, exception behavior, and complexity constraints. When applying it, preserve the distinction between mandatory requirements, recommendations, permissions, and behavior deliberately left undefined, unspecified, or implementation-defined.
\end{CornellSummaryBox}

\end{document}
